\begin{longtable}{|p{3.4cm}|p{10.3cm}|}
\caption{Đặc tả kịch bản thu hồi thẩm quyền của bên phát hành}
\label{tab:issuer-revocation-scenario} \\

\hline
\textbf{Mục} & \textbf{Nội dung} \\
\hline
\endfirsthead

\hline
\textbf{Mục} & \textbf{Nội dung} \\
\hline
\endhead

\hline
\endfoot

\hline
\endlastfoot

Mã kịch bản & KB5 \\
\hline

Tác nhân chính & Quản trị viên gốc hoặc quản trị viên \\
\hline

Thành phần tham gia &
Ví của người quản trị, mô-đun \cbox{authority}, tầng chỉ mục, dịch vụ xác minh \\
\hline

Mục đích &
Đình chỉ hoặc thu hồi thẩm quyền cấp phát của một đơn vị, sao cho các lần xác minh sau đó phát hiện được trạng thái này \\
\hline

Điều kiện tiên quyết &
Người thực hiện là quản trị viên gốc hoặc quản trị viên đang hoạt động; đơn vị bị xử lý có bản ghi thẩm quyền \\
\hline

Phân biệt ba thao tác &
\begin{itemize}[leftmargin=*, nosep]
    \item \textbf{Đình chỉ}: trạng thái chuyển sang tạm đình chỉ. Quyền vẫn còn trong bản ghi, chỉ bị khóa. Vai trò trong chỉ mục \textbf{không đổi}
    \item \textbf{Kích hoạt lại}: đưa trạng thái từ tạm đình chỉ trở về đang hoạt động
    \item \textbf{Thu hồi}: trạng thái chuyển sang đã thu hồi. Vai trò trong chỉ mục hạ xuống chủ thể, nhưng đơn vị vẫn giữ định danh và các văn bằng của riêng mình
\end{itemize} \\
\hline

Luồng thực hiện chính &
\begin{enumerate}[leftmargin=*, nosep]
    \item Người quản trị chọn đơn vị, chọn thao tác và nhập lý do
    \item Ví ký giao dịch gọi hàm tương ứng trên mô-đun \cbox{authority}
    \item Hợp đồng cập nhật trạng thái, ghi lý do và dấu thời gian, rồi phát sự kiện
    \item Tầng chỉ mục xác minh giao dịch, cập nhật vai trò nếu cần và ghi nhật ký kiểm toán
    \item Từ thời điểm này, bước thứ bảy của quy trình xác minh phát hiện trạng thái mới
\end{enumerate} \\
\hline

Hệ quả đối với các chứng thực đã cấp &
\begin{itemize}[leftmargin=*, nosep]
    \item Các chứng thực do đơn vị đó cấp \textbf{vẫn hợp lệ về mặt kỹ thuật}: chữ ký vẫn đúng, nội dung vẫn khớp giá trị băm đã neo
    \item Nhưng chúng \textbf{bị từ chối ở bước kiểm thẩm quyền}, và báo cáo xác minh nêu rõ nguyên nhân là thẩm quyền của bên phát hành chứ không phải lỗi nội dung
    \item Đơn vị \textbf{vẫn thu hồi được} các chứng thực mình đã cấp, theo nguyên tắc bất đối xứng ở Mục 3.4.5
\end{itemize} \\
\hline

Luồng ngoại lệ &
\begin{itemize}[leftmargin=*, nosep]
    \item Người thực hiện không có quyền $\rightarrow$ hợp đồng từ chối
    \item Kích hoạt lại một đơn vị không ở trạng thái tạm đình chỉ $\rightarrow$ từ chối với mã lỗi tương ứng
    \item Thu hồi một đơn vị đã bị thu hồi $\rightarrow$ từ chối, tránh ghi hai bản ghi cho cùng một thao tác
\end{itemize} \\
\hline

Kết quả sau cùng &
Trạng thái mới tra cứu được công khai trên chuỗi; bản ghi cũ giữ nguyên làm dấu vết; thao tác có mã giao dịch và một dòng nhật ký \\
\hline

Tính chất được kiểm chứng &
F3, F10, F12, N1, N2 --- và nguyên tắc bất đối xứng ở Mục 3.4.5 \\
\hline

\end{longtable}
